\documentclass[12pt]{article}
\usepackage{xeCJK}
\usepackage{fontspec}
\usepackage{amsmath}
\usepackage{graphicx}
\usepackage{geometry}
\usepackage{siunitx}
\usepackage{float}
\usepackage{fancyhdr}
\usepackage{circuitikz}
\usepackage{caption}
\usepackage{subcaption}
\usepackage{lastpage}
\usepackage{booktabs}

\captionsetup[figure]{
    font=footnotesize,
    name=图
}
\captionsetup[table]{
    font=footnotesize,
    name=表
}
\setCJKmainfont{Songti SC}
\geometry{a4paper, margin=1in, top=1in}
\setlength{\parindent}{2em}

% 页眉页脚设置
\pagestyle{fancy}
\fancyhf{}
\setlength{\headheight}{14pt}
\fancyhead[R]{\small 普通物理实验报告}
\fancyhead[C]{\small 实验名称：虚拟仪器实验}
\fancyfoot[C]{\small 第 \thepage\ 页/共 \pageref{LastPage} 页}
\fancypagestyle{plain}{%
    \fancyhf{}
    \fancyfoot[C]{\small 第 \thepage\ 页/共 \pageref{LastPage} 页}
    \renewcommand{\headrulewidth}{0pt}  % 移除页眉分割线
}

\title{\Large\textbf{虚拟仪器实验报告}}
\author{普通物理实验2\\
         \\
        姓名：\underline{\hspace{1cm}[REDACTED]\hspace{1cm}} \\
        学号：\underline{\hspace{0.5cm}[REDACTED]\hspace{0.5cm}} \\
        序号：\underline{\hspace{1.4cm}[REDACTED]\hspace{1.4cm}}}
\date{\today}

\begin{document}

\maketitle
\section{第一周：虚拟仪器基础测量}

\subsection{电阻测量 ($50\,\Omega$ 与 $1k\,\Omega$)}

    \par
    
    \hspace{2em}利用搭建的虚拟仪器程序对标称值为 $50\,\Omega$ 和 $1\,\text{k}\Omega$ 的电阻进行了测量。

\begin{figure}[H]
    \centering
    \caption{$50\,\Omega$、$1k\,\Omega$ 电阻伏安特性测量前面板截图}
    \begin{subfigure}{0.48\textwidth}
        \centering
        \includegraphics[width=0.95\textwidth]{report_images/50ohm.jpg}
    \end{subfigure}
    \hfill
    \begin{subfigure}{0.48\textwidth}
        \centering
        \includegraphics[width=1.05\textwidth]{report_images/1kohm.jpg}
    \end{subfigure}
\end{figure}

测量数据记录如下表所示：

\begin{table}[H]
    \centering
    \caption{电阻测量数据记录表}
    \begin{tabular}{ccccc}
        \toprule
        \textbf{电阻标称值} & 测量值 1 ($\Omega$) & 测量值 2 ($\Omega$) & 测量值 3 ($\Omega$) & \textbf{平均值 ($\Omega$)} \\
        \midrule
        $50\,\Omega$ & 49.8381 & 49.8339 & 49.8465 & 49.8395 \\
        $1\,\text{k}\Omega$ & 1088.63 & 1089.21 & 1089.37 & 1089.07 \\
        \bottomrule
    \end{tabular}
\end{table}

\textbf{测量不确定度评定：}

本次测量仅考虑测量数据的重复性引入的误差，即\textbf{A类测量不确定度}（Type A Uncertainty）。根据贝塞尔公式计算实验标准偏差 $s$，则平均值的标准不确定度 $u_A = s / \sqrt{n}$（其中 $n=3$）。

\begin{itemize}
    \item 对于 $50\,\Omega$ 电阻，计算标准不确定度
    \[
        \sigma_{50} = \sqrt{\frac{\sum (R_{50,i} - \bar{R_{50}})^2}{3(3-1)}} \approx 0.0037\,\Omega
    \]
    
    测量结果表示为：$R_{50} = 49.840 \pm 0.004\,\Omega$
    
    \item 对于 $1\,\text{k}\Omega$ 电阻，计算标准不确定度
    \[
        \sigma_{1k} = \sqrt{\frac{\sum (R_{1k,i} - \bar{R_{1k}})^2}{3(3 - 1)}} \approx 0.22\,\Omega
    \]
    测量结果表示为：$R_{1k} = 1089.1 \pm 0.2\,\Omega$
\end{itemize}

\subsection{二极管伏安曲线测量}

\hspace{2em}利用虚拟仪器测量了二极管的正向和反向伏安特性曲线。

\begin{figure}[H]
    \centering
    \begin{subfigure}{0.53\textwidth}
        \centering
        \includegraphics[width=1.0\textwidth]{report_images/UI diodode.jpg}
        \caption{二极管伏安特性测量结果}
    \end{subfigure}
    \hfill
    \begin{subfigure}{0.45\textwidth}
        \centering
        \includegraphics[width=\textwidth]{data/week 01/UI diodode diagram2.jpg}
        \caption{二极管伏安特性测量程序框图}
    \end{subfigure}
\end{figure}

\textbf{主要修改内容：}
\begin{itemize}
    \item 修改终止条件：电流扫描范围在 $\pm 8\,\text{mA}$ 之间 （绝对值小于$8\,\text{mA}$）。
    \item 实现正反向扫描：引入一个新的shift register，用于控制电压的扫描，每次对电压加上步长重新放入shift register，替代原本的循环坐标i。
    \item 删除线性拟合部分相关代码。
    \item 调整步长为$0.1V$，扫描范围$-6$ 至 $2V$，充分展现二极管伏安特性。
\end{itemize}
\textbf{(选做) 静态电阻计算：}

\par
\vspace{1em}
根据测量结果，在电流 $I = \pm 4\,\text{mA}$ 处计算二极管的静态电阻 $R = U/I$。
由于测量数据点较为密集（$\Delta I \approx 0.1\,\text{mA}$），在极小的电流变化范围内二极管伏安特性可视为线性变化，因此采用\textbf{线性插值法}计算目标电流对应的电压值。
\vspace{0.5em}
\begin{table}[H]
    \centering
    \caption{二极管静态电阻测量数据}
    \begin{tabular}{cc|cc}
        \toprule
        \textbf{电流1 (mA)} & \textbf{电压1 (V)} & \textbf{电流2 (mA)} & \textbf{电压2 (V)} \\
        \midrule
        -4.0017 & -5.3578 & -3.9151 & -5.3566 \\
        3.9375 & 0.77615 &  4.0315 & 0.77692 \\
        \bottomrule
    \end{tabular}
\end{table}
\vspace{0.5em}
\par
根据选取的 $I = \pm 4\,\text{mA}$ 附近的两个数据点进行插值运算：
\begin{itemize}
    \item 正向 4mA 处（利用 $3.9375\,\text{mA}$ 与 $4.0315\,\text{mA}$ 两点插值）：\\
    $U \approx$ 0.77667  V, \quad $R_{+} = U / 4\,\text{mA} \approx$ \underline{ 194.17 } $\Omega$。
    
    \item 反向 -4mA 处（利用 $-4.0010\,\text{mA}$ 与 $-3.9151\,\text{mA}$ 两点插值）：\\
    $U \approx$ -5.3578  V, \quad $R_{-} = U / (-4\,\text{mA}) \approx$ \underline{ 1339.5 } $\Omega$。
\end{itemize}

\newpage

\section{第二周：Fano 共振实验}

\subsection{实验电路与原理}

\hspace{2em}本实验旨在研究 Fano 共振现象。Fano 共振是一种量子干涉效应，表现为分立态与连续态之间的干涉，导致光谱呈现不对称的线型。在宏观电路中，可以通过耦合振荡回路来模拟这一现象。 \\

\par

实验电路主要由信号源、耦合电容 $C$ 以及并联谐振回路（由电感 $L2$、谐振电容 $C_2$ 和损耗电阻 $R_1$ 组成）构成。

\begin{figure}[h]
    \centering
    \begin{circuitikz}[american, scale=1.0, transform shape]
        % --------------------------------------------------
        % 1. 定义坐标节点
        % --------------------------------------------------
        \coordinate (GND) at (0,4);   % 左上角 AI GND
        \coordinate (AI0) at (0,0);   % 左下角 AI 0
        \coordinate (TopJunction) at (8,4); % 上方节点 (C与R0连接处)
        \coordinate (BotJunction) at (8,0); % 下方节点 (C与C1连接处)
        
        % --------------------------------------------------
        % 2. 绘制左侧电源与R0回路
        % --------------------------------------------------
        % 电源 Us (注意: 实际连接在 AI GND 和 AI 0 之间)
        \draw (GND) node[ground]{} node[above]{AI GND}
            to[sV, l=$U_S$] (AI0) node[below]{AI 0};
            
        % 顶部电阻 R0
        \draw (GND) to[R, l=$R_0$, a=$100\,\Omega$] (TopJunction)
            node[above ] {AI 1}; % 红色标注 AI 1 点

        % --------------------------------------------------
        % 3. 绘制底部主回路 (R1 -> L1 -> C1)
        % --------------------------------------------------
        \draw (AI0) 
            to[R, l=$R_1$ ] (2.5,0)
            to[L, l=$L_1$, a=$18\,\text{mH}$] (5.0,0)
            to[C, l=$C_1$, a=$0.047\,\mu\text{F}$] (BotJunction);

        % --------------------------------------------------
        % 4. 绘制共用耦合元件 (C)
        % --------------------------------------------------
        % 垂直放置的电容 C
        \draw (BotJunction) 
            to[C, l=$C$] (TopJunction);

        % --------------------------------------------------
        % 5. 绘制右侧回路 (C2 -> L2)
        % --------------------------------------------------
        % 从下方节点引出 -> C2 -> L2 -> 回到上方节点
        \draw (BotJunction) -- (9,0)
            to[C, l=$C_2$] (11,0)
            to[L, l=$L_2$, a=$16\,\text{mH}$] (13,0)
            -- (13,4) % 向上折线
            -- (TopJunction);

    \end{circuitikz}
    \caption{共用元件实现耦合的 RLC 测量电路图}
\end{figure}

\section{电路幅频特性理论推导与参数分析}

\subsection{幅频特性表达式推导}
\hspace{2em}
为得到归一化幅频特性（即导纳 $Y(\omega) = I/U$），我们将电路划分为两个主要阻抗部分：

\begin{enumerate}
    \item \textbf{背景串联支路阻抗 $Z_1$}（对应连续态）：由 $R_1, L_1, C_1$ 串联而成。
    \begin{equation}
        Z_1(\omega) = R_1 + j\left(\omega L_1 - \frac{1}{\omega C_1}\right)
    \end{equation}
    
    \item \textbf{共振并联模块阻抗 $Z_2$}（对应离散态与耦合）：由耦合电容 $C$ 与谐振支路 $(L_2, C_2)$ 并联而成。
    \begin{itemize}
        \item 谐振支路阻抗：$Z_{LC} = j\left(\omega L_2 - \frac{1}{\omega C_2}\right)$
        \item 耦合电容阻抗：$Z_C = \frac{1}{j\omega C}$
    \end{itemize}
    根据并联电路公式 $Z_2 = \frac{Z_C \cdot Z_{LC}}{Z_C + Z_{LC}}$，代入整理得：
    \begin{equation}
        Z_2(\omega) = \frac{j\left(\omega L_2 - \frac{1}{\omega C_2}\right)}{1 - \omega C \left(\omega L_2 - \frac{1}{\omega C_2}\right)}
    \end{equation}
\end{enumerate}

电路的总负载阻抗为 $Z_{total} = Z_1 + Z_2$。因此，归一化幅频特性（导纳模值）的严格表达式为：

\begin{equation}
    \left| \frac{I}{U} \right| = \left| Y(\omega) \right| = \frac{1}{\left| Z_1(\omega) + Z_2(\omega) \right|} 
    = \frac{1}{\left| R_1 + j(\omega L_1 - \frac{1}{\omega C_1}) + \frac{j(\omega L_2 - \frac{1}{\omega C_2})}{1 - \omega C (\omega L_2 - \frac{1}{\omega C_2})} \right|}
    \label{eq:strict_admittance}
\end{equation}

\subsection{电路参量对 Fano 线型的影响}
\hspace{2em}
从式 (\ref{eq:strict_admittance}) 可以直接分析各电路参量如何控制 Fano 线型（波峰与波谷）的形态：

\begin{itemize}
    \item \textbf{$C_2$ 决定波峰位置：}
    当 $Z_{LC} \approx 0$ 即 $\omega L_2 = 1/(\omega C_2)$ 时，并联模块被“短路”，阻抗 $Z_2 \to 0$，总阻抗最小，电流达到最大值。
    \begin{equation}
        \omega_{\text{peak}} \approx \frac{1}{\sqrt{L_2 C_2}}
    \end{equation}
    \textbf{结论：} 改变 $C_2$ 会直接平移 Fano 形状中的\textbf{峰值频率}，而对波形形状影响较小。

    \item \textbf{$C$ 决定波谷位置及不对称度：}
    当 $Z_2$ 的分母趋近于 0，即发生并联谐振时，模块阻抗 $Z_2 \to \infty$，电流趋近于 0，形成波谷。条件为 $1 = \omega C (\omega L_2 - 1/\omega C_2)$。
    \begin{equation}
        \omega_{\text{dip}} \approx \sqrt{\frac{1}{L_2}(\frac{1}{C_2} + \frac{1}{C})}
    \end{equation}
    \textbf{结论：} 改变耦合电容 $C$ 会改变波谷的位置。
    \begin{itemize}
        \item 波峰与波谷的频率差 $\Delta \omega = |\omega_{\text{dip}} - \omega_{\text{peak}}|$ 决定了曲线的“陡峭程度”和非对称因子 $q$。
        \item $C$ 越小，$\Delta \omega$ 越大，波谷离波峰越远，耦合效应越弱。
        \item $C$ 越大，波谷向波峰靠近，耦合效应越强。
    \end{itemize}

    \item \textbf{$R_1$ 决定背景高度与尖锐度：}
    $R_1$ 存在于分母的实部。
    \textbf{结论：} 
    \begin{itemize}
        \item $R_1$ 增大，整体导纳 $|Y|$ 下降，背景基线降低。
        \item $R_1$ 同时也提供了阻尼，共振峰变得更加平滑（Q值降低），甚至掩盖掉细微的共振特征。
    \end{itemize}
\end{itemize}

\subsection{元件表征值测量}
\hspace{2em}
实验首先对 Fano 共振电路中使用的主要元件（电感 $L_1$、$L_2$ 和电容 $C_1$）进行了频率特性表征。测量结果如下图所示：

\begin{figure}[H]
    \centering
    \begin{subfigure}[b]{0.48\textwidth}
        \centering
        \includegraphics[width=\textwidth]{report_images/component_16mH.png}
        \caption{L=16mH 电感测量结果}
    \end{subfigure}
    \hfill
    \begin{subfigure}[b]{0.48\textwidth}
        \centering
        \includegraphics[width=\textwidth]{report_images/component_18mH.png}
        \caption{L=18mH 电感测量结果}
    \end{subfigure}
    \caption{电感元件表征}
\end{figure}

\begin{figure}[H]
    \centering
    \includegraphics[width=0.6\textwidth]{report_images/component_047uF.png}
    \caption{C=0.047$\mu$F 电容测量结果}
\end{figure}

\subsection{Fano 共振幅频特性曲线的复现}
\hspace{2em}使用标准值0.2$\mu$F的谐振电容$C_2$，0.5$\mu$F的耦合电容$C$，以及500$\Omega$的损耗电阻$R$，测量了Fano共振电路的幅频特性曲线。实验测量曲线如下图：
\begin{figure}[H]
    \caption{Fano 共振标准值幅频特性曲线复现}
    \centering
    \includegraphics[width=0.7\textwidth]{report_images/fano_standard.png}
\end{figure}


\subsection{Fano 共振幅频特性曲线研究}
\hspace{2em}
通过改变 Fano 共振电路中的参数（耦合电容 $C$、谐振电容 $C_2$、损耗电阻 $R$），测量了多组幅频特性曲线。为了清晰展示参数变化对共振曲线的影响，我们将同一元件不同参数下的曲线绘制在同一张图中。图中\textbf{深色实线}为 B-Spline 拟合曲线，\textbf{浅色散点}为实验测量数据。

\subsubsection{改变谐振电容 $C_2$ 的影响}

保持耦合电容 $C = 0.5\,\mu\text{F}$ 和 $R = 500\,\Omega$ 不变，改变谐振回路电容 $C_2$ 的值。

\begin{figure}[H]
    \centering
    \includegraphics[width=0.9\textwidth]{report_images/fano_comparison_vary_c2.png}
    \caption{不同谐振电容 $C_2$ 下的 Fano 共振幅频曲线}
\end{figure}

\textbf{分析与讨论：}
$C_2$ 是并联谐振回路的核心元件之一。根据前述分析，谐振频率 $f_0 \approx \frac{1}{2\pi\sqrt{L_2C_2}}$。
从图中可以清晰地观察到，随着 $C_2$ 的增大（从 $0.03\,\mu\text{F}$ 到 $0.3\,\mu\text{F}$），共振峰的位置显著向低频方向移动。这与公式预测完全一致。
同时，共振峰的宽度和高度也随之改变，这反映了电路品质因数 $Q$ 的变化。在 $C_2$ 较小时，共振频率较高，电路的特征阻抗较大，导致更尖锐细小的共振峰。

\subsubsection{改变耦合电容 $C$ 的影响}
\hspace{2em}
保持谐振电容 $C_2 = 0.2\,\mu\text{F}$ 和 $R = 500\,\Omega$ 不变，改变耦合电容 $C$ 的值。

\begin{figure}[H]
    \centering
    \includegraphics[width=0.9\textwidth]{report_images/fano_comparison_vary_c.png}
    \caption{不同耦合电容 $C$ 下的 Fano 共振幅频曲线}
\end{figure}

\textbf{分析与讨论：}
电容 $C$ 在电路中起耦合作用，连接信号源与谐振回路。它决定了“连续态”（背景信号）与“分立态”（谐振回路信号）之间的干涉强度。
从图中可以看出，随着 $C$ 的变化，共振曲线的形状发生显著变化，不对称性程度发生变化：
\begin{itemize}
    \item 较小的 $C$ 意味着弱耦合，背景通道信号较弱，共振主要由 $LC$ 回路主导，曲线更接近对称的洛伦兹型。
    \item 较大的 $C$ 增强了直通信号，使得干涉效应增强，导致线型出现明显的非对称结构。
\end{itemize}

\subsubsection{改变电阻 $R$ 的影响}
\hspace{2em}
保持 $C = 0.5\,\mu\text{F}$ 和 $C_2 = 0.2\,\mu\text{F}$ 不变，改变电阻 $R$ 的值。

\begin{figure}[H]
    \centering
    \includegraphics[width=0.9\textwidth]{report_images/fano_comparison_vary_r.png}
    \caption{不同电阻 $R$ 下的 Fano 共振幅频曲线}
\end{figure}

\textbf{分析与讨论：}
从图中可以看出，随着电阻 $R$ 从 $5\,\Omega$ 增加到 $10\,\text{k}\Omega$，Fano 共振的特征峰谷逐渐变浅、变宽。
    \begin{itemize}
        \item 当 $R$ 较小（如 $5\,\Omega$, $50\,\Omega$）时，曲线表现出极高的锐度，波谷较深，非对称性显著。
        \item 当 $R$ 增大至 $10\,\text{k}\Omega$时，共振特征几乎完全消失，幅频曲线退化为一条平坦的直线。
    \end{itemize}
    
    \textbf{分析：}电阻 $R$ 在电路中主要引入\textbf{阻尼效应}。
    \begin{itemize}
        \item 小电阻对应弱阻尼，连续态背景的能量损耗小，与离散态的干涉效应强烈，导致显著的 Fano 振荡。
        \item 大电阻引入强阻尼，使得背景共振峰过度展宽，最终掩盖了精细的Fano共振特征。
    \end{itemize}

\section{总结}
\hspace{2em}
本次实验通过虚拟仪器成功复现了 Fano 共振现象，并系统研究了电路参数对共振线型的影响。实验结果表明：
实验曲线与 Fano 理论公式定性吻合，验证了该电路模型在模拟量子干涉Fano共振效应方面的有效性。
\end{document}
